% Example of the MVP Lab technical report template (mvplab.sty).
% Compile with pdfLaTeX:  latexmk -pdf example.tex
% (or: pdflatex example && bibtex example && pdflatex example && pdflatex example)
\documentclass[10pt,letterpaper]{article}
\usepackage{mvplab}

% Optional: packages of your own go here; mvplab.sty loads hyperref and cleveref at the end of the preamble.
\usepackage{multirow}

% ---------------------------------------------------------------------------------------------------------------
% Title data
\title{An Example Technical Report Typeset with the MVP Lab Template}
\shorttitle{An Example Technical Report with the MVP Lab Template} % running header (optional)
% \projectlead marks the project leader (blue dagger), \corresponding the corresponding author (pink envelope).
% Write names with ~ so that a line break never splits a name.
\author{First~Author\projectlead \and Second~Author \and Third~Author \and Fourth~Author\corresponding}
\affiliation{MVP Lab \& Imperial College London}
% Several affiliations: put \affilmark{1}, \affilmark{2}, ... after the names and before each affiliation:
%   \author{First~Author\affilmark{1}\projectlead \and Second~Author\affilmark{2}}
%   \affiliation{\affilmark{1}MVP Lab, Imperial College London\qquad\affilmark{2}Another Institution}
\mvpdate{\today}
% Every link line is optional; leave a command out to drop its line.
\projecturl{https://mvp-ai-lab.github.io/your-project}
\githuburl{https://github.com/mvp-ai-lab/your-project}
\huggingfaceurl{https://huggingface.co/datasets/mvp-lab/your-dataset}
\projectleader{First Author}{first.author@example.com}
\correspondence{Fourth Author}{fourth.author@example.com}

\begin{document}

% The abstract comes before \maketitle: it is typeset inside the title box, above the date and the links.
\begin{abstract}
This document shows the MVP Lab template for technical reports and arXiv papers. The first page carries the lab
logo, the title, the authors with the project leader and the corresponding author marked, and the abstract in a box
together with the date, the project page, the code and data repositories, and the contact addresses. Every later
page has the title as running header. Citations are numbered~\cite{dust3r,vggt}, and citations as well as references
to sections, figures, tables, and equations are colored in the accent pink of the lab's webpages (\cref{sec:usage}).
The remaining sections show headings, a figure, a table, equations, lists, boxes, and an appendix.
\end{abstract}
\maketitle

% ---------------------------------------------------------------------------------------------------------------
\section{Introduction}
\label{sec:intro}
The template is a single \texttt{.sty} file for the standard \texttt{article} class. It sets a single-column letter
page with one-inch margins, Latin Modern for the text and Inter for headings, captions, and the title page, and the
colors of the MVP Lab webpages: primary blue, indigo, and accent pink. Links to web pages are blue, \eg
\url{https://github.com/mvp-ai-lab}, and internal references are pink, \eg \cref{fig:example,tab:example} and
\cref{eq:loss}.

Recent geometric foundation models predict depth and camera poses from images in a single forward
pass~\cite{dust3r,vggt}, while monocular depth models predict depth from one image~\cite{midas,depthanything,marigold}.
We use them only as example citations; the bibliography style is numbered and sorted, and consecutive numbers are
compressed into ranges~\cite{midas,depthanything,marigold,dust3r,vggt}.

\paragraph{Contributions.} A run-in paragraph heading is set in bold in the ink color:
\begin{itemize}
  \item a title page with logo, author roles, links, and a boxed abstract;
  \item headings, captions, and links in the colors of the lab's webpages;
  \item a running header with the title on every page after the first, and an appendix on a new page.
\end{itemize}

% ---------------------------------------------------------------------------------------------------------------
\section{Using the Template}
\label{sec:usage}

\subsection{Title Page}
\label{sec:usage:title}
The title data are set in the preamble (see the top of \texttt{example.tex}): \verb|\title|, an optional
\verb|\shorttitle| for the running header, \verb|\author| with \verb|\and| between names, \verb|\affiliation|, and the
optional \verb|\projecturl|, \verb|\githuburl|, \verb|\huggingfaceurl|, \verb|\projectleader|, \verb|\correspondence|,
and \verb|\mvpdate|. The abstract environment comes before \verb|\maketitle|.

\subsection{Figures and Tables}
\label{sec:usage:floats}
Figures and tables use the standard environments; captions get a blue label, and table captions go above the table.
\Cref{fig:example} is drawn with TikZ in the template colors, and \cref{tab:example} uses \texttt{booktabs}.

\begin{figure}[t]
\centering
\begin{tikzpicture}[x=1.15cm,y=0.045cm]
  \draw[mvpmuted!60] (0,0) -- (8.6,0);
  \foreach \y in {0,25,50,75,100} {
    \draw[mvpbg2] (0,\y) -- (8.6,\y);
    \node[left,font=\sffamily\scriptsize,text=mvpmuted] at (0,\y) {\y};
  }
  \foreach \x/\a/\b/\name in {1/42/58/Task A,3/55/71/Task B,5/63/80/Task C,7/48/77/Task D} {
    \fill[mvpblue] (\x-0.38,0) rectangle (\x,\a);
    \fill[mvppink] (\x,0) rectangle (\x+0.38,\b);
    \node[below,font=\sffamily\scriptsize,text=mvpink] at (\x,0) {\name};
  }
  \node[rotate=90,font=\sffamily\scriptsize,text=mvpmuted] at (-0.75,50) {Score};
  \fill[mvpblue] (5.9,104) rectangle (6.1,110); \node[right,font=\sffamily\scriptsize] at (6.1,107) {Baseline};
  \fill[mvppink] (7.2,104) rectangle (7.4,110); \node[right,font=\sffamily\scriptsize] at (7.4,107) {Ours};
\end{tikzpicture}
\caption{\textbf{An example figure.} Scores of a baseline and of our method on four tasks (made-up numbers).}
\label{fig:example}
\end{figure}

\begin{table}[t]
\centering
\caption{\textbf{An example table.} Lower is better for AbsRel, higher for $\delta_{1.25}$; the best result of each
column is in bold (made-up numbers).}
\label{tab:example}
\small
\begin{tabular}{@{}llcccc@{}}
\toprule
& & \multicolumn{2}{c}{Indoor} & \multicolumn{2}{c}{Outdoor} \\
\cmidrule(lr){3-4}\cmidrule(l){5-6}
Method & Type & AbsRel $\downarrow$ & $\delta_{1.25}$ $\uparrow$ & AbsRel $\downarrow$ & $\delta_{1.25}$ $\uparrow$ \\
\midrule
\multirow{2}{*}{Baselines} & Monocular & 0.081 & 93.2 & 0.102 & 89.5 \\
 & Multi-view & 0.074 & 94.1 & 0.095 & 90.8 \\
\midrule
Ours & Multi-view & \textbf{0.062} & \textbf{95.7} & \textbf{0.083} & \textbf{92.6} \\
\bottomrule
\end{tabular}
\end{table}

\subsection{Equations}
\label{sec:usage:eq}
Equations are centered and numbered on the right. A single equation:
\begin{equation}
  \mathcal{L}(\theta) = \frac{1}{N}\sum_{i=1}^{N} \rho_\delta\!\left(\frac{a\,\hat{x}_i(\theta)+b-x_i}{m}\right),
  \label{eq:loss}
\end{equation}
and several equations aligned at the equals sign:
\begin{align}
  \mathrm{AbsRel} &= \frac{1}{|\Omega|}\sum_{p\in\Omega}\frac{|\tilde{D}(p)-D(p)|}{D(p)}, \label{eq:absrel}\\
  \delta_{\tau} &= \frac{1}{|\Omega|}\sum_{p\in\Omega}\mathbb{1}\!\left[\max\!\left(\frac{\tilde{D}(p)}{D(p)},
    \frac{D(p)}{\tilde{D}(p)}\right)<\tau\right]. \label{eq:delta}
\end{align}

\subsection{Boxes}
\label{sec:usage:boxes}
The template defines a rounded box, \verb|mvpbox|, which takes any \texttt{tcolorbox} option, for example a title:
\begin{mvpbox}[title=Key finding]
Boxes highlight a finding, a definition, or a summary. Their border, background, and title follow the template
colors.
\end{mvpbox}
The \verb|mvpprompt| variant sets its text in small typewriter type, for prompts and listings (\cref{sec:appx:prompt}).

% ---------------------------------------------------------------------------------------------------------------
\section{Conclusion}
\label{sec:conclusion}
To start a report, copy \texttt{mvplab.sty} and the \texttt{template/} folder (the logos) next to your main
file. Then load the package after \verb|\documentclass{article}| and fill in the title data as in this example.

{
  \small
  \bibliographystyle{plainnat}
  \bibliography{example}
}

% ---------------------------------------------------------------------------------------------------------------
% Appendix: a new page with its own heading and no running header; sections are lettered.
\appendixpage
\appendix

\section{Prompt Box}
\label{sec:appx:prompt}
A prompt in the \verb|mvpprompt| box:
\begin{mvpprompt}
The image shows an indoor scene. Estimate the metric depth of each numbered point, in meters, and return a single
JSON object mapping each point number (as a string) to its depth, wrapped in <answer> </answer> tags.
\end{mvpprompt}

\section{Further Results}
\label{sec:appx:results}
Appendix sections are lettered (\cref{sec:appx:prompt,sec:appx:results}), and figures, tables, and equations continue
the numbering of the main text.

\end{document}
